\chapter{分段、索引、保留与截断}
单文件日志随时间增长，会令启动扫描越来越慢，也无法单独删除旧数据。Kafka 风格的日志用多个 segment 组成一个分区：旧 segment 封闭，当前段承接追加；稀疏索引把逻辑 offset 映射到某个段内字节边界。索引是可重建的加速结构，不是记录真相。日志保留和截断是不同操作：前者删除旧的封闭前缀，后者把尾部回退到某个位点以修复分叉。

\begin{figure}[ht]
\centering
\begin{tikzpicture}[x=0.8mm,y=1mm,font=\small]
\draw[thick,rounded corners] (0,0) rectangle (57,22); \node at (28.5,16) {segment base 000};
\node at (28.5,8) {offset 0--3};
\draw[thick,rounded corners] (65,0) rectangle (122,22); \node at (93.5,16) {segment base 004};
\node at (93.5,8) {offset 4--7};
\draw[thick,rounded corners,fill=blue!5] (130,0) rectangle (187,22); \node at (158.5,16) {active base 008};
\node at (158.5,8) {LEO=8};
\draw[-{Stealth},thick] (28, -4) -- (28, -13); \node[below] at (28,-13) {index floor(2)=(0,pos)};
\draw[-{Stealth},thick] (94, -4) -- (94, -13); \node[below] at (94,-13) {index floor(6)=(4,pos)};
\end{tikzpicture}
\caption{分段日志的逻辑范围与段内稀疏索引。段文件名中的 base offset 是该段逻辑起点。}
\end{figure}

\lessonsection{5}{稀疏索引}
\goals{构建可恢复的 offset 到 byte-position 稀疏映射。需要理解 floor 查找、记录边界与日志扫描恢复。}
\background{索引项格式是两个大端 int64：\pathcode{offset, position}，每项 16 字节。每隔 intervalRecords 条记录写一个点，第一条记录总要有索引项。查找目标 offset 时，floor 返回不大于目标的最大 offset，使读取可从近处合法边界开始并继续顺扫。不存在合适项时，\pathcode{(baseOffset,0)} 是安全起点。索引可能丢失、截短或含非法项；主日志仍是事实来源，恢复时应扫描日志重建索引。}
\providededit{索引文件路径、段 baseOffset、interval 和生命周期已在 \method{SparseIndex} 提供；\method{SegmentLog} 可从记录重扫。}{\pathcode{exercises/src/main/java/io/simplekafka/storage/SparseIndex.java}：\method{add(long,long)}、\method{floor(long)}、\method{rebuild(Path,long)}。}
\lessoncontract{添加项按 offset 严格递增且 position 非负；floor 返回最大 offset<=目标，找不到时返回段起点的零位置。索引缺失、文件长度不是 16 的倍数或条目非法时从日志重建；日志本身损坏仍须失败。重建不能把错误索引伪装成可信日志数据。}
\lessonhints{先写清楚 floor 的数学定义：候选集合是哪些 offset，空集合的约定是什么。}{把索引读取失败分成“索引可丢弃”和“主日志不可恢复”两类。}{用多段记录构造间隔索引，删除索引文件后再通过真实读取观察恢复效果，而不只检查内存 map。}
\expected{interval=4 时，对 offset 0、4、8 建项；floor(6) 返回 offset 4 对应的物理位置。删去 \pathcode{.index} 后 reopen 并读取 offset 6，仍得到正确记录。非法长度索引文件也可重建；有损坏日志仍报错。对应 \method{Step05Test}。}
\pitfalls{索引项若指向记录内部，读取会把任意 payload 当作 header；floor 误用 ceiling 会跳过目标前的边界。索引若成为唯一数据来源，缺索引就丢日志。Kafka 的索引组织与精度/字节偏移模型更复杂；课程每 intervalRecords 条记录索引一次，不是按固定字节间隔。}
\lessoncommand{5}
\lessonquestions{为什么找到 floor entry 后仍要扫描记录？}{损坏的 index 和损坏的 log 为何不能采用同一种恢复策略？}

\lessonsection{6}{分段日志}
\goals{把单段日志组合成可滚动的分区日志，并让读取跨 segment 连续工作。需要理解段边界、容量估算、索引与全局 LEO。}
\background{PartitionLog 按 baseOffset 排序持有多个 SegmentLog。活动段若容纳不了下一条合法记录，就封闭当前段并新建 base=当前 LEO 的段；单条记录即使大于配置 segmentBytes，只要记录自身合法，也应独占新段而非永远无法追加。一次 batch 可以横跨段边界，offset 在整个 partition 范围连续。读取先定位目标段，然后用该段索引 floor 得到提示，顺扫后如有需要进入下一段，并把 maxRecords/maxBytes 作为跨段总预算。创建新文件前需要预检整个 batch 的合法性，避免已知非法输入改变目录结构。}
\providededit{每个 segment 的 create/open/recover、稀疏索引与资源关闭都已提供；\method{PartitionLog} 维护按 baseOffset 排序的段表。}{\pathcode{exercises/src/main/java/io/simplekafka/storage/PartitionLog.java}：\method{append(List<RecordData>)}、\method{read(long,int,int)}。}
\lessoncontract{append 返回整个 batch 的 [first,next) 区间，跨段 offset 仍无间隙。容量不足时只在追加前滚段；非法 batch 不得先滚段。读操作按 offset 顺序跨段，累计总记录与字节限制；段直接追加本身不负责滚动。重开已存在目录后恢复 LEO，不重复分配旧 offset。目录发现只处理约定 \pathcode{.log/.index} 文件，不删除陌生文件。}
\lessonhints{容量判断应该使用编码后的完整记录字节数，而不是只计算 value 长度。}{决定一次 batch 的每条记录与段边界如何配合，同时保持全局 offset 连续。}{将测试分成初次 append、滚动后 append、关闭重开后的 append，再逐条核对所有 offset。}
\expected{配置 256 字节段并追加十条记录，关闭重开再追加一条，读取范围应仍为 0 到 10 且没有重复 offset。测试也覆盖 batch 正好跨段、单条记录独占新段、非法超长记录前后目录/LEO不变；对应 \method{Step06Test}。}
\pitfalls{容量只看原始 payload 会漏算固定头、长度和 CRC；只让每个 batch 单独检查 offset 会造成跨段冲突。真实 Kafka segment 更常按字节规模滚动并结合索引文件和时间策略；课程也用容量滚动，但索引间隔是记录数。单段 append 不自动滚段是有意分层。}
\lessoncommand{6}
\lessonquestions{为何要在创建任何新段之前验证整批记录？}{跨段 fetch 的字节预算应如何避免在段切换时被重置？}

\lessonsection{7}{日志保留}
\goals{安全删除已过期的封闭段，同时保留 active segment 与当前 LEO。需要区分日志起点、消费进度和单条删除。}
\background{分区日志实际只有一段可保留范围 \pathcode{[logStartOffset, LEO)}。课程用显式调用进行保留，不运行后台定时器。\method{deleteBefore(offset)} 只删除最后记录 offset 小于阈值的封闭段，不能按记录切割一段，也不能移除活动段；新起点必须落在真实段边界。消费者若请求已经删除的数据会得到 \method{OFFSET_OUT_OF_RANGE}，不静默把 position 改到新起点。消费组 offset 提交与物理删日志独立，推进消费位点不会自动触发保留。}
\providededit{多段目录扫描、日志起点和段关闭路径由 \method{PartitionLog} 骨架持有；步骤 6 的段表与索引可以复用。}{\pathcode{exercises/src/main/java/io/simplekafka/storage/PartitionLog.java}：\method{deleteBefore(long)}。}
\lessoncontract{只删满足条件的封闭段及同 base 的 \pathcode{.index}，不删 active，不删除未知文件。返回实际新的 logStartOffset。LEO 不倒退；重开后新起点仍正确。若已有读取位点低于新起点，读取必须明确报范围错误，不自动 seek。}
\lessonhints{从段 base 与下一段 base 确认哪些逻辑 offset 属于可删前缀。}{空 active 段与有数据 active 段都要单独考虑，不能把“最后一个段”当成可删项。}{删除后关闭并重新打开，再从边界前后分别读取，以验证磁盘文件和内存起点一致。}
\expected{构造三个段、保留最后一个活动段并删除旧封闭段后，实际 logStartOffset 移到保留边界；重开仍返回相同起点且 LEO 不变。低于新起点读取报 \method{OFFSET_OUT_OF_RANGE}。对应 \method{Step07Test}。}
\pitfalls{只删 \pathcode{.log} 不删配对索引会制造悬空索引；按阈值删掉 active 会令后续追加失去起点。真实 Kafka 保留常由时间/大小配置与后台管理触发，并可能有 retention 与 consumer 独立关系；课程只提供明确同步调用，不替消费者重置位点。}
\lessoncommand{7}
\lessonquestions{为什么删除阈值通常不会恰好对应某条记录的 segment 边界？}{提交位点已经越过旧段，是否意味着系统可以自动删段？为什么？}

\lessonsection{8}{尾部截断}
\goals{把日志尾部截回某个下一条位点，删除分叉尾数据并恢复可追加状态。需要理解 half-open 区间、记录边界和段 base。}
\background{截断目标 \pathcode{nextOffset} 表示保留所有 offset<nextOffset 的记录；合法范围为 [logStartOffset,LEO]。若目标恰在已有段边界，应有可追加的 active segment 以该 base 开始；若位于某段内部，应截断到该记录的物理边界，删除后续 segment，并重建被保留段索引。截到 LEO 是无变化；截到 logStartOffset 留下空段，下一条 append 使用这个 offset。截断只改变本地日志尾，不会自动改变复制状态、HW 或消费位点。}
\providededit{段关闭、文件截断、索引重建和全局段表已由前面步骤建立；资源管理仍由 PartitionLog 提供。}{\pathcode{exercises/src/main/java/io/simplekafka/storage/PartitionLog.java}：\method{truncateTo(long)}。}
\lessoncontract{只允许 [logStartOffset,LEO] 内的 nextOffset，超范围拒绝且日志不变。保留 offset<nextOffset，所有之后数据及其索引移除；LEO 更新到 nextOffset。成功之后可立刻追加且新记录的 offset 正好等于 nextOffset。未知文件不可删除。}
\lessonhints{将目标理解成“下一条要写的 offset”，而不是要保留的最后一条 offset。}{分别列出截在段中、截在段边界、截在当前 LEO、截在 logStartOffset 的磁盘状态。}{重开截断后的目录并追加一条；只检查内存 nextOffset 不足以证明索引和文件同步。}
\expected{写入 offset 0–8 后执行 \pathcode{truncateTo(5)}，offset 0–4 仍可读；已删除的 5–8 不会返回，读取新的 LEO 5 得到空结果，下一条追加得到 offset 5。\pathcode{truncateTo(LEO)} 不改变内容；截到起点后保留空 active 并从该起点追加。对应 \method{Step08Test}。}
\pitfalls{误将 nextOffset 解释为最后保留记录，会产生 off-by-one；只 truncate 文件不重建索引，可能使之后读取跳到删除的分叉记录。真实 Kafka 截断往往与 leader epoch、HW 和副本协议联动；本步只提供日志原语，不允许拿它单独推断安全选举。}
\lessoncommand{8}
\lessonquestions{截断到 5 时，为何下一条应写入 offset 5 而非 4？}{为什么清理文件尾之后还必须处理索引文件？}
